\expandafter\def\csname HistoricalNoteText24\endcsname{Die Juno erreicht am Ende des Aprils 1807 ihre Sonnenferne.}
\expandafter\def\csname HistoricalNoteText25\endcsname{Ich theilte hierauf Hrn. Dr. \textsc{Gauss} meine Beobachtungen der Ceres am M. Q. vom 10. und 12. Jan. mit, so wie noch einige von Hrn. Dr. \textsc{Koch} in Danzig angestellten, die nachher folgen. \textsc{Bode.}}
\expandafter\def\csname HistoricalNoteText26\endcsname{Die auf der Mailänder Sternwarte angestellte und von Herrn \textsc{Oriani} uns gütigst überschickte Beobachtung der Opposition dieses Planeten, welche wir hier zunächst folgen lassen, war zu jener Zeit, als sich Hr. Dr. \textsc{Gauss} mit Berechnung der neuen Elemente beschäftigte, wegen des durch die Kriegsumstände verursachten unordentlichen Postenlaufs noch nicht angelangt. Der Brief, welcher sie enthielt, war beinahe sechs Wochen unterwegs. [\textsc{von Zach.}]}
\clearpage
\fancyhead[L]{\small GAUSS AN BODE.}
\hypertarget{native-283-lower}{}
\begin{center}
GAUSS AN BODE.
\end{center}

\begin{center}
Astronomisches Jahrbuch für 1809. S. 215–219. Berlin 1806.
\end{center}

\begin{flushright}
Braunschweig, den 30. Juli 1806.
\end{flushright}

\par\indent In meinem letzten Briefe habe ich Ew. --- das Versprechen gegeben, Ihnen die Resultate meiner
neuesten Rechnungen über die neuen Planeten mitzutheilen, sobald ich damit fertig sein würde; ich eile
jetzt dasselbe zu erfüllen.

\par\indent Meine im Februar d. J. gemachten Beobachtungen der Pallas habe ich Ihnen bereits angezeigt
(S. Seite 137 [S. 267 d. B.]). Zu der letzten Verbesserung der Elemente habe ich mich indess nur der zu
Mailand gemachten Beobachtungen bedient, die nun auch im Juli-Heft der M. C. abgedruckt sind. Dies
sind demnach meine neuesten elliptischen Elemente der Pallas.

\begin{center}
\emph{Elemente der Pallas.} (IX.)

Epochen der Länge, Meridian von Seeberg.
\end{center}

\begin{center}\begingroup
\fontsize{9}{11}\selectfont
\setlength{\tabcolsep}{3pt}\renewcommand{\arraystretch}{1.12}
\begin{tabular}{ll}
1802.  & $143^\circ\,24'\,1''1$ \\
1803.  & $221\;32\;56.4$ \\
1804.  & $299\;54\;42.5$ \\
1805.  & $18\;3\;37.8$ \\
1806.  & $96\;12\;33.2$ \\
1807.  & $174\;21\;28.5$ \\
\end{tabular}\endgroup\end{center}
\begin{center}\begingroup
\fontsize{9}{11}\selectfont
\setlength{\tabcolsep}{3pt}\renewcommand{\arraystretch}{1.12}
\begin{tabular}{ll}
tägliche mittlere Bewegung .\ .\ .\ .\ . & $770''7816$ \\
Sonnenferne 1806 .\ .\ .\ .\ .\ .\ .\ . & $301^\circ\,12'\,21''$ \\
aufsteigender Knoten 1806 .\ .\ .\ . & $172\;30\;56$ \\
Neigung der Bahn .\ .\ .\ .\ .\ .\ . & $34\;37\;8$ \\
Excentricität .\ .\ .\ .\ .\ .\ .\ .\ .\ . & $0.245384$ \\
Logarithm der halben Axe .\ .\ .\ .\ . & $0.4421021$ \\
\end{tabular}\endgroup\end{center}
\SourceLeaf{284}%
\par\indent Nach diesen Elementen habe ich den Lauf der Pallas für ihre nächste Erscheinung berechnet,
welche Ephemeride ich hier beizufügen die Ehre habe. Prof. \textsc{Harding} hat abermals eine Karte dazu
verfertigt, die Hr. \textsc{v. Zach} für die M. C. stechen lassen wird. Von Prof. \textsc{Harding}'s Atlas, der für die
Beobachtung der Asteroiden so sehr nothwendig ist, werden nächstens schon ein paar Blätter erscheinen.

\par\indent Auch von der Juno habe ich Ihnen die hier gemachten Beobachtungen bereits mitgetheilt (S.
Seite 137 [S. 267 d. B.]); so viel ich bis jetzt sehe, waren es die ersten in diesem Jahre. Auch hier habe
ich zur Verbesserung der Elemente die \textsc{Oriani}'schen Beobachtungen, die im Juli-Hefte der M. C. abgedruckt sind, hauptsächlich gebraucht.

\begin{center}
VI. \emph{Elemente der Juno.}

Epochen der mittlern Länge. Meridian von Seeberg.
\end{center}

\begin{center}\begingroup
\fontsize{9}{11}\selectfont
\setlength{\tabcolsep}{3pt}\renewcommand{\arraystretch}{1.12}
\begin{tabular}{ll}
1805.  & $42^\circ\,35'\,7''3.$ \\
1806.  & $125^\circ\,11'\,20''1.$ \\
1807.  & $207^\circ\,47'\,33''0$ \\
\end{tabular}\endgroup\end{center}

\begin{center}\begingroup
\fontsize{9}{11}\selectfont
\setlength{\tabcolsep}{3pt}\renewcommand{\arraystretch}{1.12}
\begin{tabular}{ll}
Länge der Sonnenferne 1806 .\ .\ . & $233^\circ\,17'\,1''$ \\
Länge des aufsteigenden Knotens 1806  & $171\;4\;57$ \\
Neigung der Bahn .\ .\ .\ .\ .\ .\ . & $13\;3\;28''4$ \\
tägliche mittl. tropische Bewegung . & $814''720$ \\
Excentricität .\ .\ .\ .\ .\ .\ .\ .\ . & $0.2549441$ \\
Logarithm der halben Axe .\ .\ .\ . & $0.4260480$ \\
\end{tabular}\endgroup\end{center}

\par\indent Die nach diesen Elementen berechnete Ephemeride habe ich auch schon Hrn. Prof. \textsc{Harding} mitgetheilt, um danach die Karte zu zeichnen: Sie erhalten hiebei eine Abschrift davon. Es wird im nächsten Jahre aber äusserst schwer halten, die Juno zu beobachten: schon in diesem Jahre war sie äusserst
lichtschwach, und im nächsten wird sie nicht einmal halb so hell erscheinen als in diesem\HistoricalNote{24}. [Der Lauf
der Juno vom 7. Dec. 1806 bis 21. Sept. 1807 ist mit dreitägigen Zwischenräumen für 12 Uhr Seeberg bis
auf Minuten der Geraden Aufsteigung und der Abweichung genau angegeben.] Die $\text{\AstroSymbols ☍}$ der $\PallasSymbol$ trifft hienach ein 1807 den 4. Mai um 21 Uhr im $223^\circ\,53'$ Länge und $42^\circ\,18'$ nördl. Breite. 

\par\indent Über die Ceres habe ich dieses Jahr noch keine weitere Rechnungen angestellt; es fehlt mir noch
an guten Beobachtungen in der Nähe der Opposition. Da die $\CeresSymbol$ eine so günstige Lage hatte, so zweifle
ich nicht, dass diese $\text{\AstroSymbols ☍}$ an mehreren Orten gut beobachtet sein wird. Sollten Sie im Besitz von
dergleichen Beobachtungen sein, so erzeigen Sie mir eine Gefälligkeit, wenn Sie mir dieselben \emph{bald} mittheilen. Gern werde ich Ihnen dann auch die Resultate davon zuschicken\renewcommand{\thefootnote}{**)}\HistoricalNote{25}\renewcommand{\thefootnote}{*)}.



\Needspace{15\baselineskip}
\fancyhead[L]{\small GAUSS AN VON ZACH.}
\begin{center}
GAUSS AN VON ZACH.
\end{center}

\begin{center}
Monatliche Correspondenz. Band XIV. S. 377–382. 1806. October.
\end{center}

\begin{flushright}
Braunschweig, den 25. August 1806.
\end{flushright}

\par\indent --- --- --- Beigehend erhalten Sie die Ephemeride für die Juno. Sie ist zwar schon seit einiger Zeit fertig gewesen, allein immer hoffte ich, dass ich zugleich auch die für die \emph{Ceres} würde mit
beifügen können, von der ich aber immer noch keine ganz zuverlässigen Beobachtungen aus der Nähe
\SourceLeaf{285}%
der Opposition habe. --- Sollte ich nicht bald gute Beobachtungen derselben erhalten, so werde ich
doch einstweilen die Verbesserung der Elemente mit den Beobachtungen von \textsc{Pasquich}, die nach seinem
eignen Urtheil nicht sehr zuverlässig sind, machen müssen. \textsc{Harding}'s und ein paar mir von \textsc{Bode} mitgetheilte Beobachtungen liegen schon etwas weiter von der Opposition ab. Sollte dieselbe nicht auch
in Mailand beobachtet sein?

\par\indent Die Elemente, wonach die Ephemeride der \emph{Juno} berechnet ist, sind folgende:

\begin{center}
VI. \emph{Elemente der Juno.}
\end{center}

\Needspace{18\baselineskip}
\begin{center}
Epoche der Länge im Meridian von Seeberg
\end{center}
\[
\left\{\begin{array}{ll}
1805. & 42^\circ\,35'\,7''3\\
1806. & 125\;11\;20.1\\
1807. & 207\;47\;33.0
\end{array}\right.
\]
\begin{center}\begingroup
\fontsize{9}{11}\selectfont
\setlength{\tabcolsep}{3pt}\renewcommand{\arraystretch}{1.12}
\begin{tabular}{ll}
tägliche mittlere tropische Bewegung  & $814''7201$ \\
Länge der Sonnenferne 1806 .\ .\ . & $233^\circ\,17'\,1''$ \\
Länge des aufsteigenden Knotens 1806  & $171^\circ\,4'\,57''4$ \\
Excentricität .\ .\ .\ .\ .\ .\ .\ .\ . & $0.254944$ \\
Logarithm der halben Axe .\ .\ .\ . & $0.4260480$ \\
Neigung der Bahn .\ .\ .\ .\ .\ .\ . & $13^\circ\,3'\,28''4$ \\
\end{tabular}\endgroup\end{center}

\par\indent Prof. \textsc{Harding} schreibt mir, dass er beide Karten schon grösstentheils vollendet habe, eine Unpässlichkeit verzögere nur noch die Vergleichung mit dem Himmel selbst. Er hofft sie Ihnen bald zuschicken zu können.

\Needspace{15\baselineskip}
\par\indent Die Lichtstärke der \emph{Juno} im nächsten Jahre habe ich in meiner Ephemeride nicht mit angeführt,
weil ich sie diesmal nach einer neuen ganz vorzüglich bequemen Methode berechnet habe, bei der aber
der Abstand des Planeten von der Sonne in der Rechnung gar nicht vorkommt. Zur Übersicht setze ich
aber doch drei berechnete Grössen für die Lichtstärke her, woraus erhellt, wie ausserordentlich schwer
die Beobachtung der Juno im nächsten Jahre sein wird:

\begin{center}\begingroup
\fontsize{9}{11}\selectfont
\setlength{\tabcolsep}{3pt}\renewcommand{\arraystretch}{1.12}
\begin{tabular}{ll}
1807 Januar 12 .\ . & $0.00661,$ \\
Mai 18 .\ . & $0.01600,$ \\
Septbr. 21 .\ . & $0.00709$ \\
\end{tabular}\endgroup\end{center}

\par\indent In diesem Jahre war im März die grösste Lichtstärke doch noch $0.03983$ (M. C. XI. B. S. 478. [S. 265
d. B.]), also 1807, wenn sie am grössten ist, nicht einmal halb so gross. --- Die Abkürzung der Rechnung bei der eben gedachten neuen Methode ist so gross, dass ich auf eine gewöhnliche Octavseite 15
vollständige Rechnungen bringe, wenn die Sonnen-Örter als schon gegeben angesehen werden.

\par\indent [Der geocentrische Lauf der Juno von 1807 Jan. 12 bis Sept. 21 ist mit dreitägigen Zwischenräumen für 12 U. Mitternacht in Seeberg dargestellt, die Gerade Aufsteigung und Abweichung bis auf Minuten, der Abstand von der Erde bis auf drei Decimalen genau.]

\Needspace{34\baselineskip}
\begin{center}
GAUSS AN VON ZACH.
\end{center}

\begin{center}
Monatliche Correspondenz. Band XV. S. 152–157. 1807 Februar.
\end{center}

\begin{flushright}
Braunschweig, den 3. Januar 1807.
\end{flushright}

\par\indent --- --- --- Ich habe die Ehre, Ihnen hier die Ephemeride für den Lauf der \emph{Ceres} einzusenden
die Elemente zur Berechnung derselben hatte ich aus Ermangelung anderer\HistoricalNote{26} nur auf Herrn Professor
\SourceLeaf{286}%
\textsc{Pasquich}'s Beobachtung der letzten Opposition (M. C. XIII. B. S. 192; XIV. B. S. 92.) gründen
können, und muss mir also ihre fernere Verbesserung bis zum Empfang zuverlässigerer vorbehalten.
Es waren folgende:

\begin{center}
XI. \emph{Elemente.}
\end{center}

\begin{center}\begingroup
\fontsize{9}{11}\selectfont
\setlength{\tabcolsep}{3pt}\renewcommand{\arraystretch}{1.12}
\begin{tabular}{ll}
\multicolumn{2}{l}{Epoche der Länge im Meridian von} \\
\quad Seeberg 1806 .\ .\ .\ .\ .\ .\ .\ . & $108^\circ\,19'\,34''7$ \\
Länge der Sonnenferne 1806 .\ .\ . & $326\;37\;59$ \\
Länge des aufsteigenden Knotens$\cdot$1806  & $80\;53\;23$ \\
Neigung der Bahn .\ .\ .\ .\ .\ .\ . & $10\;37\;33.7$ \\
Mittlere jährliche tropische Bewegung  & $78^\circ\,9'\,23''3$ \\
Mittlere tägliche tropische Bewegung . & $770''8584$ \\
Excentricität .\ .\ .\ .\ .\ .\ .\ .\ .\ . & $0.0783486$ \\
Logar. der halben Achse .\ .\ .\ .\ .\ . & $0.4420728$ \\
\end{tabular}\endgroup\end{center}

\par\indent Mein Werkchen über die Bestimmung der Planeten-Bahnen u.\,s.\,w., dem ich den grössten Theil
meiner Zeit gewidmet habe, naht sich seiner Vollendung; ich habe bisher 22 Bogen Manuscript
ausgearbeitet und bin eben jetzt beschäftigt, die ausführlichen Beispiele einzutragen. Eins ist von
der \emph{Juno} hergenommen, wo die Zwischenzeit 22 Tage ist; bei dem zweiten von der \emph{Pallas} ist
sie 71 Tage; das dritte von der \emph{Ceres}, wo sie 118 Tage war, habe ich wieder cassirt, um durch ein
anderes, wo sie 260 Tage beträgt, die grosse Allgemeinheit der Methode noch besser in's Licht zu
setzen. In allen diesen Beispielen lässt sie sich unmittelbar so anwenden, dass noch gar nichts von
der Bahn als bekannt angesehen wird. Das Ganze wird etwa 30 Bogen im Druck betragen. Bei den
jetzigen Zeitumständen, fürchte ich, wird es vielleicht in Deutschland einige Mühe kosten, einen
soliden Verleger zu finden. --- --- --- [Der geocentrische Lauf der Ceres vom 19. Dec. 1806 bis
21. Sept. 1807 ist ebenso wie ich oben Seite 278 d. B. für die Pallas angedeutet habe, dargestellt,
nur ist die Lichtstärke bis auf fünf statt dort auf vier Delimalen genau bestimmt.]
